\chapter{परिमित समूहों के निरूपण}\label{ch:b3:representations}

किसी अमूर्त समूह को समझने के लिए उसे किसी सदिश समष्टि पर क्रिया
कराइए और उस क्रिया पर रैखिक बीजगणित --- अभिलक्षणिक मान, अनुरेख,
अंतःगुणन --- लगाइए। यह कार्यक्रम, अर्थात् \emph{\omterm{def:b3:representations:rep}{निरूपण} सिद्धांत},
$\C$ पर परिमित समूहों के लिए आश्चर्यजनक रूप से कारगर है: प्रत्येक
\omterm{def:b3:representations:rep}{निरूपण} \omterm{def:b3:representations:rep}{अखंडनीय} निरूपणों में टूट जाता है (मास्के), \omterm{def:b3:representations:rep}{अखंडनीय निरूपण}
अपने \emph{वर्णों} (अनुरेखों) से निर्धारित हो जाते हैं, और \omterm{def:b3:representations:character}{वर्ण}
ऐसे लांबिकता संबंध संतुष्ट करते हैं जो परिकलन को यांत्रिक बना
देते हैं। यह अध्याय इसी कलन को और उसकी पहली उत्कृष्ट कृतियों को
--- छोटे समूहों की \omterm{def:b3:representations:character}{वर्ण} सारणियों को --- खड़ा करता है, और
सप्ताहांत समस्या एक ऐसी प्रमेय काटती है जो निरे समूह सिद्धांत की
पहुँच से बहुत दूर है: बर्नसाइड की $p^aq^b$ प्रमेय। सर्वत्र $G$ एक
परिमित समूह है और सभी सदिश समष्टियाँ $\C$ पर परिमित विमा की हैं।

\section{निरूपण, मास्के, शूर}

\begin{definition}\label{def:b3:representations:rep}
$G$ का कोई \emph{निरूपण}\index{निरूपण} किसी $\C$-सदिश समष्टि
$V$ के लिए एक समाकारिता $\rho \colon G \to GL(V)$ है; $\dim V$ उसकी \emph{घात} कहलाती है।
कोई उपसमष्टि $W \subseteq V$ \emph{अपरिवर्ती} कहलाती है यदि सभी $g$ के लिए
$\rho(g)W \subseteq W$ हो; प्रतिबंध से $W$ एक \emph{उपनिरूपण} बन जाता है। $\rho$
\emph{अखंडनीय}\index{अखंडनीय निरूपण} कहलाता है यदि $V \neq
0$ हो और
उसकी एकमात्र अपरिवर्ती उपसमष्टियाँ $0$ और $V$ हों। $(\rho, V)$ और
$(\sigma, W)$ के बीच \emph{समाकारिता} ऐसा रैखिक $f\colon V \to W$ है जिसके लिए सभी
$g$ पर $f\rho(g) = \sigma(g)f$ हो (समतुल्यवर्तिता); एकैकी आच्छादक $f$
\emph{तुल्याकारिताएँ} कहलाती हैं।
\end{definition}

\begin{example}\label{ex:b3:representations:examples}
(क) घात $1$: समाकारिताएँ $G \to \C^\times$।
(ख) \emph{नियमित निरूपण}\index{नियमित निरूपण}: आधार $(e_h)_{h \in G}$ वाला
$V = \C^G$, जिसमें $\rho(g)e_h = e_{gh}$; घात $\abs G$।
(ग) किसी परिमित समुच्चय $X$ पर $G$ की क्रमचय क्रिया $\C^X$ पर
\emph{क्रमचय \omterm{def:b3:representations:rep}{निरूपण}} देती है: $\rho(g)e_x =
e_{g\cdot x}$।
(घ) $S_n$ निर्देशांकों का क्रमचय करके $\C^n$ पर क्रिया करता है;
अधिसमतल $\{\sum x_i = 0\}$ अपरिवर्ती है: यही घात $n - 1$ का \emph{मानक \omterm{def:b3:representations:rep}{निरूपण}}
है।
\end{example}

\begin{theorem}[मास्के]\label{thm:b3:representations:maschke}
किसी \omterm{def:b3:representations:rep}{निरूपण} $(V, \rho)$ की प्रत्येक अपरिवर्ती उपसमष्टि $W$ का कोई
अपरिवर्ती पूरक होता है। फलस्वरूप प्रत्येक \omterm{def:b3:representations:rep}{निरूपण} \omterm{def:b3:representations:rep}{अखंडनीय}
निरूपणों का प्रत्यक्ष योग है
(\emph{अर्धसरलता})।\index{मास्के की प्रमेय}
\end{theorem}

\begin{proof}
मान लीजिए $p \colon V \to V$ प्रतिबिंब $W$ वाला \emph{कोई भी} प्रक्षेप है
(कोई भी पूरक चुन लीजिए)। उसका समूह पर औसत लीजिए:
\[
\tilde p = \frac1{\abs G}\sum_{g \in G} \rho(g)\,p\,\rho(g)^{-1}.
\]
प्रत्येक पद $V$ को $W$ में भेजता है ($W$ अपरिवर्ती है) और
$W$ को बिंदुशः स्थिर रखता है: $w \in W$ के लिए $\rho(g)^{-1}w \in W$, $p$ उसे स्थिर
रखता है, और $\rho(g)$ उसे वापस ले आता है --- अतः $\tilde p$ पुनः $W$ पर
एक प्रक्षेप है। वह समतुल्यवर्ती भी है: $h \in G$ के लिए $\rho(h)\tilde p\rho(h)^{-1}$ उसी योग
का पुनःसूचकांकन कर देता है। अतः $\ker
\tilde p$ $W$ का अपरिवर्ती पूरक है।
योज्यांशों पर इसे दोहराने से (परिमित विमा) $V$ \omterm{def:b3:representations:rep}{अखंडनीय}
निरूपणों में विघटित हो जाता है।
\end{proof}

\begin{theorem}[शूर की प्रमेयिका]\label{thm:b3:representations:schur}
मान लीजिए $f \colon V \to W$ \emph{\omterm{def:b3:representations:rep}{अखंडनीय}} निरूपणों की समाकारिता है। तब
$f = 0$ या $f$ एक तुल्याकारिता है; और यदि $(V,\rho) = (W,\sigma)$ हो, तो किसी $\lambda \in \C$
के लिए $f = \lambda\,\mathrm{id}$। अतः $\dim\operatorname{Hom}_G(V, W)$ $V \cong W$ होने पर $1$ है और अन्यथा
$0$।\index{शूर की प्रमेयिका}
\end{theorem}

\begin{proof}
$\ker f$ और $\operatorname{im} f$ अपरिवर्ती हैं (समतुल्यवर्तिता), अतः इनमें से
प्रत्येक $0$ है या पूरा: या तो $f = 0$, या $f$ एकैकी है और उसका
प्रतिबिंब पूरा। यदि $V = W$ हो: $f$ का कोई अभिलक्षणिक मान $\lambda$
होता है ($\C$ \omterm{def:b3:galois:closure}{बीजीय रूप से संवृत}); $f - \lambda\,\mathrm{id}$ एक अनेकैकी समाकारिता
$V \to V$ है, अतः $0$। $V \cong W$ की स्थिति में विमा गणना के लिए: कोई एक
तुल्याकारिता $u$ स्थिर कीजिए; तब कोई भी समाकारिता $f$
अंतःसमाकारिता $u^{-1}f = \lambda\,\mathrm{id}$ देती है: अतः $f =
\lambda u$।
\end{proof}

\section{वर्ण और लांबिकता}

\begin{definition}\label{def:b3:representations:character}
$(V, \rho)$ का \emph{वर्ण}\index{वर्ण} $\chi_\rho(g) = \operatorname{tr}\rho(g)$ है। वह $\chi_\rho(e) = \dim V$, $\chi_\rho(hgh^{-1}) = \chi_\rho(g)$ को
संतुष्ट करता है (अनुरेख संयुग्मन के अंतर्गत अपरिवर्ती होते हैं):
अर्थात् वर्ण \emph{वर्ग फलन}\index{वर्ग फलन} हैं --- संयुग्मन
वर्गों पर अचर फलनों की समष्टि $\mathcal{CF}(G)$ के अवयव, जिस पर एर्मीतीय
अंतःगुणन
\[
\langle \varphi, \psi\rangle = \frac1{\abs G}\sum_{g \in G}
\overline{\varphi(g)}\,\psi(g).
\]
दिया गया है।
\end{definition}

\begin{proposition}\label{prop:b3:representations:charbasics}
$\rho(g)$ विकर्णनीय है और उसके अभिलक्षणिक मान इकाई के मूल हैं; $\chi_\rho(g^{-1}) = \overline{\chi_\rho(g)}$,
और $\abs{\chi_\rho(g)} \leq \chi_\rho(e)$, जिसमें समता तभी और केवल तभी होती है जब $\rho(g)$ अदिश हो।
प्रत्यक्ष योगों पर \omterm{def:b3:representations:character}{वर्ण} जुड़ जाते हैं: $\chi_{V
\oplus W} = \chi_V + \chi_W$।
\end{proposition}

\begin{proof}
$N = \abs G$ के लिए $\rho(g)^N = \mathrm{id}$ (लाग्राँज): अतः $\rho(g)$ $X^N - 1$ को नष्ट कर देता
है, और वह सरल मूलों के साथ विभाजित होता है, इसलिए वह विकर्णनीय
है और उसके अभिलक्षणिक मान $\lambda_i \in \mu_N$ हैं। तब $\chi(g^{-1}) = \sum \lambda_i^{-1} = \sum\bar\lambda_i =
\overline{\chi(g)}$, और $\abs{\chi(g)} = \abs{\sum\lambda_i}
\leq \dim V$, जिसमें
त्रिभुज असमिका में समता तभी और केवल तभी होती है जब सभी $\lambda_i$
समान हों, अर्थात् $\rho(g) = \lambda\,\mathrm{id}$। प्रत्यक्ष योगों पर योज्यता: खंडों के
अनुरेख।
\end{proof}

\begin{lemma}\label{lem:b3:representations:homtrace}
मान लीजिए $(V, \rho)$, $(W, \sigma)$ \omterm{def:b3:representations:rep}{निरूपण} हैं। $\operatorname{Hom}(W, V)$ पर औसतकारी संकारक
\[
c(A) = \frac1{\abs G}\sum_{g}\rho(g)\,A\,\sigma(g)^{-1},
\]
$\operatorname{Hom}_G(W, V)$ पर एक प्रक्षेप है, और प्रतिचित्रण $\Phi_g \colon A \mapsto \rho(g)A\sigma(g)^{-1}$ के लिए $\operatorname{tr}\Phi_g =
\chi_\rho(g)\,\overline{\chi_\sigma(g)}$।
\end{lemma}

\begin{proof}
$c(A)$ समतुल्यवर्ती है (मास्के की तरह योग का पुनःसूचकांकन कीजिए),
और $c$ समतुल्यवर्ती प्रतिचित्रणों को स्थिर रखता है (प्रत्येक पद
$A$ के बराबर है): अतः $c$ $\operatorname{Hom}_G(W, V)$ प्रतिबिंब वाला प्रक्षेप है।
अनुरेख: आधारों में $\Phi_g(A) = BAC$, जहाँ $B = \rho(g)$, $C = \sigma(g)^{-1}$; आव्यूहों के आधार
$(E_{kl})$ पर $BE_{kl}C = \sum_{m,n}
b_{mk}c_{ln}E_{mn}$, अतः $\Phi_g(E_{kl})$ में $E_{kl}$ का गुणांक $b_{kk}c_{ll}$ है: इसलिए
$\operatorname{tr}\Phi_g =
\sum_{k,l}b_{kk}c_{ll} = \operatorname{tr}(B)
\operatorname{tr}(C) = \chi_\rho(g)\chi_\sigma(g^{-1})$, और $\chi_\sigma(g^{-1}) = \overline{\chi_\sigma(g)}$।
\end{proof}

\begin{theorem}[प्रथम लांबिकता संबंध]
\label{thm:b3:representations:orthogonality}
मान लीजिए $\rho, \sigma$ \emph{\omterm{def:b3:representations:rep}{अखंडनीय}} हैं। तब\index{वर्णों की लांबिकता}
\[
\langle\chi_\sigma, \chi_\rho\rangle =
\begin{cases}
1 & \text{यदि } \rho \cong \sigma,\\
0 & \text{अन्यथा:}
\end{cases}
\]
अर्थात् \omterm{def:b3:representations:rep}{अखंडनीय} \omterm{def:b3:representations:character}{वर्ण} $\mathcal{CF}(G)$ में एक प्रसामान्य लांबिक कुल बनाते
हैं।
\end{theorem}

\begin{proof}
किसी प्रक्षेप का अनुरेख उसके प्रतिबिंब की विमा होता है:
\[
\dim\operatorname{Hom}_G(W, V) = \operatorname{tr} c
= \frac1{\abs G}\sum_g \operatorname{tr}\Phi_g
= \frac1{\abs G}\sum_g \chi_\rho(g)\overline{\chi_\sigma(g)}
= \langle \chi_\sigma, \chi_\rho\rangle,
\]
और शूर की प्रमेयिका बाईं ओर का मान $\delta_{\rho \cong
\sigma}$ निकाल देती है।
\end{proof}

\begin{corollary}\label{cor:b3:representations:multiplicity}
$V \cong \bigoplus_i V_i^{\oplus m_i}$ को भिन्न-भिन्न \omterm{def:b3:representations:rep}{अखंडनीय} निरूपणों में विघटित कीजिए ($V_i \not\cong V_j$)।
तब $m_i = \langle
\chi_{V_i}, \chi_V\rangle$: अर्थात् बहुलताएँ --- और इसलिए तुल्याकारिता तक \omterm{def:b3:representations:rep}{निरूपण}
--- \omterm{def:b3:representations:character}{वर्ण} से निर्धारित हो जाती हैं। इसके अतिरिक्त $\langle \chi_V, \chi_V\rangle = \sum_i
m_i^2$; विशेष रूप
से $V$ \omterm{def:b3:representations:rep}{अखंडनीय} है तभी और केवल तभी जब $\langle\chi_V,
\chi_V\rangle = 1$।
\end{corollary}

\begin{proof}
$\chi_V = \sum m_i\chi_{V_i}$ (\cref{prop:b3:representations:charbasics}); प्रत्येक $\chi_{V_i}$ के साथ अंतःगुणन लीजिए और प्रसामान्य
लांबिकता का उपयोग कीजिए। समान \omterm{def:b3:representations:character}{वर्ण} वाले दो निरूपणों की बहुलताएँ
समान होती हैं, अतः वे तुल्याकारी हैं।
\end{proof}

\begin{theorem}[नियमित निरूपण]
\label{thm:b3:representations:regular}
मान लीजिए $\chi_1, \dots, \chi_r$ भिन्न-भिन्न \omterm{def:b3:representations:rep}{अखंडनीय} \omterm{def:b3:representations:character}{वर्ण} हैं, जिनकी घातें $n_i = \chi_i(e)$
हैं। \omterm{ex:b3:representations:examples}{नियमित निरूपण} बहुलताओं $m_i = n_i$ के साथ विघटित होता है;
फलस्वरूप
\[
\sum_{i=1}^{r} n_i^2 = \abs G,
\qquad
\sum_i n_i\chi_i(g) = 0 \quad (g \neq e).
\]
\end{theorem}

\begin{proof}
नियमित \omterm{def:b3:representations:character}{वर्ण}: $\chi_{\mathrm{reg}}(g) =
\#\{h : gh = h\}$, जो $g = e$ के लिए $\abs G$ है और अन्यथा $0$।
अतः $m_i = \langle \chi_i,\chi_{\mathrm{reg}}\rangle =
\frac1{\abs G}\overline{\chi_i(e)}\,\abs G = n_i$। $\chi_{\mathrm{reg}} = \sum n_i\chi_i$ का $e$ पर और $g \neq e$ पर मान लेने से दोनों
प्रदर्शित सर्वसमिकाएँ मिल जाती हैं।
\end{proof}

\begin{theorem}\label{thm:b3:representations:numberirr}
\omterm{def:b3:representations:rep}{अखंडनीय} \omterm{def:b3:representations:character}{वर्ण} $\mathcal{CF}(G)$ का एक प्रसामान्य लांबिक \emph{आधार} बनाते
हैं: अर्थात् \omterm{def:b3:representations:rep}{अखंडनीय} निरूपणों की संख्या $G$ के संयुग्मन वर्गों
की संख्या के बराबर है।
\end{theorem}

\begin{proof}
केवल पूर्णता शेष है: मान लीजिए $f \in \mathcal{CF}(G)$ प्रत्येक $\chi_i$ के लांबिक
है; हम $f = 0$ दिखाते हैं। किसी \omterm{def:b3:representations:rep}{निरूपण} $(V,\rho)$ के लिए $T_{f,\rho} = \frac{1}{\abs G}
\sum_g \overline{f(g)}\,\rho(g)$ रखिए।
वह समतुल्यवर्ती है: $h \in
G$ के लिए
\[
\rho(h)T_{f,\rho}\rho(h)^{-1} = \frac1{\abs G}\sum_g
\overline{f(g)}\rho(hgh^{-1})
= \frac1{\abs G}\sum_{g'}\overline{f(h^{-1}g'h)}\rho(g') =
T_{f,\rho}
\]
($f$ एक \omterm{def:b3:representations:character}{वर्ग फलन} है)। यदि $\rho$ घात $n$ का \omterm{def:b3:representations:rep}{अखंडनीय} हो, तो
शूर से $T_{f,\rho} = \lambda\,\mathrm{id}$ मिलता है, जहाँ
\[
\lambda = \frac{\operatorname{tr}T_{f,\rho}}{n}
= \frac{1}{n\abs G}\sum_g \overline{f(g)}\,\chi_\rho(g)
= \frac1n\,\langle f, \chi_\rho\rangle = 0 .
\]
अतः प्रत्येक \omterm{def:b3:representations:rep}{अखंडनीय} पर $T_{f,\rho} = 0$, और इसलिए (प्रत्यक्ष योग)
\emph{प्रत्येक} \omterm{def:b3:representations:rep}{निरूपण} पर --- विशेष रूप से \omterm{ex:b3:representations:examples}{नियमित निरूपण} पर।
उसे आधार सदिश $e_e$ पर लगाइए: $0 = T_{f,\mathrm{reg}}e_e =
\frac1{\abs G}\sum_g\overline{f(g)}e_g$, जिससे प्रत्येक $\overline{f(g)} = 0$
अनिवार्य हो जाता है। अतः प्रसामान्य लांबिक कुल $(\chi_i)$ $\mathcal{CF}(G)$ को
फैलाता है, जिसकी विमा संयुग्मन वर्गों की संख्या है।
\end{proof}

\begin{corollary}[स्तंभ लांबिकता]
\label{cor:b3:representations:column}
$g, h \in G$ के लिए:
\[
\sum_{i=1}^{r}\overline{\chi_i(g)}\,\chi_i(h) =
\begin{cases}
\abs{Z_G(g)} & \text{यदि } g, h \text{ संयुग्मी हैं},\\
0 & \text{अन्यथा।}
\end{cases}
\]
\end{corollary}

\begin{proof}
मान लीजिए $g_1, \dots, g_r$ वर्गों के प्रतिनिधि हैं और $c_j$ वर्गों के आकार।
$r \times r$ आव्यूह $U_{ij} =
\sqrt{c_j/\abs G}\;\chi_i(g_j)$ की पंक्तियाँ प्रसामान्य लांबिक हैं (\cref{thm:b3:representations:orthogonality} को
वर्गशः लिखने पर: $\sum_j \frac{c_j}{\abs G}\chi_i(g_j)\overline{\chi_{i'}(g_j)} =
\delta_{ii'}$), अर्थात् $UU^* = I$; और $UU^* =
I$ वाले किसी वर्ग
आव्यूह के लिए $U^*U = I$ भी होता है: अतः स्तंभ प्रसामान्य लांबिक हैं,
जिसे खोलने पर प्रदर्शित सर्वसमिका मिलती है ($\abs G/c_j = \abs{Z_G(g_j)}$, अर्थात्
संयुग्मन के लिए कक्षा--स्थायीकारी संबंध)।
\end{proof}

\begin{proposition}[एक-विमीय वर्ण; उन्नयन]
\label{prop:b3:representations:onedim}
(क) $G$ आबेली है तभी और केवल तभी जब उसके सभी \omterm{def:b3:representations:rep}{अखंडनीय} निरूपणों
की घात $1$ हो; और किसी भी $G$ के घात-$1$ वर्णों की संख्या
$[G : D(G)]$ है (वे आबेलीकरण के \omterm{def:b3:representations:character}{वर्ण} हैं)।
(ख) यदि $N \trianglelefteq G$ हो, तो $G/N$ के \omterm{def:b3:representations:rep}{अखंडनीय निरूपण} ($G \to G/N$ के साथ
संयोजन करके) ठीक $G$ के उन \omterm{def:b3:representations:rep}{अखंडनीय} निरूपणों तक उठ जाते हैं
जिनकी अष्टि $N$ को अंतर्विष्ट करती है।
\end{proposition}

\begin{proof}
(क) यदि $G$ आबेली हो, तो प्रत्येक वर्ग एकल-अवयवी है: $r = \abs G$,
और $\sum n_i^2 = \abs G$ से सभी $n_i = 1$ अनिवार्य हो जाते हैं; विलोमतः यदि सभी
$n_i = 1$ हों, तो \omterm{ex:b3:representations:examples}{नियमित निरूपण} एक-विमीय निरूपणों का योग है, अतः
$\rho_{\mathrm{reg}}(G)$ एक साथ विकर्णनीय है, इसलिए क्रमविनिमेय, और $\rho_{\mathrm{reg}}$ निष्ठ
है: अतः $G$ आबेली। घात-$1$ \omterm{def:b3:representations:rep}{निरूपण} आबेली लक्ष्य वाली
समाकारिताएँ $G \to \C^\times$ हैं: वे $G^{\mathrm{ab}} = G/D(G)$ से होकर गुजरती हैं (\cref{exo:b3:groups:9}), और
आबेली $G^{\mathrm{ab}}$ के भिन्न-भिन्न वर्णों की संख्या $\abs{G^{\mathrm{ab}}}$ है (उसके उतने
ही वर्ग हैं, और सभी घातें $1$)। (ख) प्रक्षेप के साथ संयोजन
अखंडनीयता को सुरक्षित रखता है (अपरिवर्ती उपसमष्टियाँ आपस में मेल
खाती हैं), और $N$ पर तुच्छ कोई \omterm{def:b3:representations:rep}{निरूपण} भागफल से होकर गुजरता है
(\cref{thm:b3:groups:firstiso})।
\end{proof}

\section{वर्ण सारणियाँ}

\begin{definition}\label{def:b3:representations:table}
$G$ की \emph{वर्ण सारणी}\index{वर्ण सारणी} $r \times r$ आव्यूह $\bigl(\chi_i(g_j)\bigr)$
है: पंक्तियाँ \omterm{def:b3:representations:rep}{अखंडनीय} वर्णों से और स्तंभ संयुग्मन वर्गों से
सूचकित हैं (और उनके आकार भी दर्शाए जाते हैं)। पंक्तियाँ भारित
गुणन के लिए प्रसामान्य लांबिक हैं और स्तंभ लांबिक (\cref{cor:b3:representations:column}):
सारणी अत्यंत अधिनिर्धारित है, और यही उसे परिकलनीय बनाता है।
\end{definition}

\begin{example}[$S_3$ की सारणी]\label{ex:b3:representations:s3}
वर्ग: $e$ (आकार 1), स्थानांतरण (3), $3$-चक्र (2); अतः $r
= 3$
\omterm{def:b3:representations:rep}{अखंडनीय}, जिनकी घातें $\sum n_i^2 = 6$ के साथ $n_i$ हैं: $1, 1,
2$। घात $1$:
तुच्छ $\mathbf 1$ और चिह्न $\varepsilon$। अंतिम पंक्ति स्तंभ लांबिकता से (अथवा
$\chi_{\mathrm{std}} = \chi_{\mathrm{perm}} - \mathbf 1$ से) निकलती है:
\[
\begin{array}{c|ccc}
S_3 & e & (1\,2)\ [3] & (1\,2\,3)\ [2]\\
\hline
\mathbf 1 & 1 & 1 & 1\\
\varepsilon & 1 & -1 & 1\\
\chi_{\mathrm{std}} & 2 & 0 & -1
\end{array}
\]
जाँच: $\langle\chi_{\mathrm{std}},\chi_{\mathrm{std}}\rangle =
\frac{1}{6}(4 + 0 + 2) = 1$: \omterm{def:b3:representations:rep}{अखंडनीय}।
\end{example}

\begin{example}[$S_4$ की सारणी]\label{ex:b3:representations:s4}
वर्ग: $e$ [1], स्थानांतरण [6], द्विक स्थानांतरण [3], $3$-चक्र
[8], $4$-चक्र [6]: पाँच \omterm{def:b3:representations:rep}{अखंडनीय}, $\sum n_i^2 =
24$, जिनमें दो घात-$1$ के
हैं ($\mathbf 1, \varepsilon$; $[S_4 :
D(S_4)] = [S_4 : A_4] = 2$): अतः घातें $1, 1, 2, 3, 3$। घात-$2$ वाला $S_4/V \cong S_3$ से उठता
है (\cref{prop:b3:representations:onedim}(ख), जहाँ $V$ क्लाइन समूह है); और घात $3$: मानक
\omterm{def:b3:representations:rep}{निरूपण} तथा $\varepsilon$ से उसका मोड़:
\[
\begin{array}{c|ccccc}
S_4 & e\,[1] & (1\,2)\,[6] & (1\,2)(3\,4)\,[3] &
(1\,2\,3)\,[8] & (1\,2\,3\,4)\,[6]\\
\hline
\mathbf 1 & 1 & 1 & 1 & 1 & 1\\
\varepsilon & 1 & -1 & 1 & 1 & -1\\
\chi_2 & 2 & 0 & 2 & -1 & 0\\
\chi_{\mathrm{std}} & 3 & 1 & -1 & 0 & -1\\
\varepsilon\chi_{\mathrm{std}} & 3 & -1 & -1 & 0 & 1
\end{array}
\]
($\chi_{\mathrm{std}}(g) = \operatorname{fix}(g) - 1$; $\chi_2$ प्रत्येक वर्ग के $V$ के सापेक्ष प्रतिबिंब पर
$S_3$-सारणी का मान लेता है।) सभी पंक्ति और स्तंभ लांबिकता जाँचें
सफल होती हैं --- उनमें से दो चलाना \cref{exo:b3:representations:3} का प्रारंभिक अभ्यास है।
\end{example}

\begin{method}\label{met:b3:representations:table}
\omterm{def:b3:representations:table}{वर्ण सारणी} बनाने के लिए: (1) \omterm{ex:b3:groups:actions}{संयुग्मन वर्ग} और उनके आकार सूचीबद्ध
कीजिए; (2) $G/D(G)$ से घात-$1$ वर्णों की संख्या गिनकर उन्हें
लिखिए; (3) $\sum n_i^2 = \abs G$ से शेष घातें ज्ञात कीजिए (छोटे पूर्णांकों की
संचयिकी); (4) सस्ते \omterm{def:b3:representations:rep}{अखंडनीय} प्राप्त कीजिए: भागफलों से उठाइए,
क्रमचय वर्णों में से $\mathbf 1$ घटाइए ($\langle\chi,\chi\rangle = 1$ जाँचिए), और ज्ञात वर्णों
को घात-$1$ वर्णों से गुणा कीजिए; (5) अज्ञात पंक्तियाँ स्तंभ
लांबिकता से पूरी कीजिए --- प्रत्येक स्तंभ पहले से पूर्ण स्तंभों
के लांबिक होता है, और स्तंभ $e$ घातें धारण करता है। सब कुछ पूरी
लांबिकता जाँच से सत्यापित कीजिए।
\end{method}

\begin{omfigure}
\begin{tikzpicture}[scale=0.52]
  \draw[thick] (0,0) rectangle (6,6);
  \fill[omDef!25] (0,5) rectangle (1,6);
  \fill[omProp!25] (1,4) rectangle (2,5);
  \fill[omThm!20] (2,0) rectangle (6,4);
  \draw (0,5) rectangle (1,6);
  \draw (1,4) rectangle (2,5);
  \draw (2,0) rectangle (6,4);
  \node[font=\small] at (0.5,5.5) {$1$};
  \node[font=\small] at (1.5,4.5) {$1$};
  \node[font=\small] at (4,2) {$M_2(\C)$};
\end{tikzpicture}

{\small $S_3$ का \omterm{ex:b3:representations:examples}{नियमित निरूपण}, खंड-विकर्ण रूप में: $\C[S_3] \cong \C \times \C \times M_2(\C)$, विमाएँ
$1 + 1 +
4 = 6 = \abs{S_3}$। सामान्य स्थिति में $\C[G] \cong \prod_i
M_{n_i}(\C)$: सर्वसमिका $\sum n_i^2 = \abs G$ आव्यूह खंडों के
बारे में एक कथन है।}
\end{omfigure}

\section{अभ्यास}

\begin{exercise}[$\star$]\label{exo:b3:representations:1}
(क) दर्शाइए कि $\Z/n\Z$ के \omterm{def:b3:representations:rep}{अखंडनीय} \omterm{def:b3:representations:character}{वर्ण} $\chi_k(\bar m) = \eu^{2\iu\pi km/n}$, $k = 0, \dots, n-1$, हैं, और $\Z/4\Z$
की \omterm{def:b3:representations:table}{वर्ण सारणी} लिखिए।
(ख) उस पर दोनों लांबिकता संबंध सत्यापित कीजिए --- और आव्यूह को
पहचानिए: इस पुस्तक में यह पहले कहाँ आ चुका है?
\end{exercise}

\begin{exercise}[$\star$]\label{exo:b3:representations:2}
मान लीजिए $G$ किसी परिमित समुच्चय $X$ पर क्रिया करता है और
$\chi$ क्रमचय \omterm{def:b3:representations:rep}{निरूपण} $\C^X$ का \omterm{def:b3:representations:character}{वर्ण} है।
(क) दर्शाइए कि $\chi(g) = \abs{\operatorname{Fix}_X(g)}$ और $\langle \mathbf 1, \chi\rangle = \#\{\text{कक्षाएँ}\}$ --- अर्थात् बर्नसाइड की गणना
प्रमेयिका (\cref{exo:b3:groups:5}) एक \omterm{def:b3:representations:character}{वर्ण} परिकलन है।
(ख) मान लीजिए क्रिया संक्रमणीय है, अतः $\chi = \mathbf 1 +
\psi$। दर्शाइए कि $\psi$
\omterm{def:b3:representations:rep}{अखंडनीय} है तभी और केवल तभी जब क्रिया \emph{$2$-संक्रमणीय} हो
(भिन्न बिंदुओं के क्रमित युग्मों पर संक्रमणीय)। \emph{($\langle\chi,\chi\rangle$ को
$X \times X$ पर कक्षाओं की संख्या के रूप में परिकलित कीजिए।)}
(ग) निष्कर्ष निकालिए कि $S_n$ ($n \geq
2$) का मानक \omterm{def:b3:representations:rep}{निरूपण} \omterm{def:b3:representations:rep}{अखंडनीय} है।
\end{exercise}

\begin{exercise}[$\star$]\label{exo:b3:representations:3}
\cref{met:b3:representations:table} का अनुसरण करते हुए $S_3$ की सारणी आरंभ से बनाइए, फिर \cref{ex:b3:representations:s4}
की $S_4$ सारणी में दो पंक्ति और दो स्तंभ लांबिकता संबंध सत्यापित
कीजिए। $\{1,2,3,4\}$ पर क्रिया करते $S_4$ के क्रमचय \omterm{def:b3:representations:character}{वर्ण} को और \omterm{def:b3:representations:character}{वर्ण}
$\chi_{\mathrm{std}}^2$ (बिंदुशः वर्ग) को \omterm{def:b3:representations:rep}{अखंडनीय} वर्णों में विघटित कीजिए।
\end{exercise}

\begin{exercise}[$\star\star$]\label{exo:b3:representations:4}
$D_4$ और $Q_8$ की \omterm{def:b3:representations:character}{वर्ण} सारणियाँ परिकलित कीजिए। निष्कर्ष निकालिए
कि दो अतुल्याकारी समूहों की \omterm{def:b3:representations:character}{वर्ण} सारणियाँ समान हो सकती हैं ---
फिर भी इस युग्म में सारणी कौन-सा समूह-सैद्धांतिक आँकड़ा पकड़
लेती है (केंद्रों के क्रम, आबेलीकरण, स्वप्रतिलोमी अवयवों की
संख्या)? और इनमें से किसे वह पकड़ने में \emph{विफल} रहती है?
\end{exercise}

\begin{exercise}[$\star\star$]\label{exo:b3:representations:5}
$A_4$ की \omterm{def:b3:representations:table}{वर्ण सारणी}: वर्ग $e$ [1], द्विक स्थानांतरण [3], तथा
$3$-चक्रों के \emph{दो} वर्ग [4], [4]।
(क) $3$-चक्रों के विभाजन को समझाइए (\cref{exo:b3:groups:11} की तरह $S_4$ और
$A_4$ में केंद्रकों की तुलना कीजिए)।
(ख) तीनों घात-$1$ \omterm{def:b3:representations:character}{वर्ण} ($A_4/V \cong
\Z/3\Z$ के माध्यम से) और घात-$3$ \omterm{def:b3:representations:character}{वर्ण}
($S_4$ से $\chi_{\mathrm{std}}$ को प्रतिबंधित कीजिए) ज्ञात कीजिए, और सारणी
संकलित कीजिए।
(ग) सारणी से $A_4$ के \omterm{def:b3:groups:normal}{प्रसामान्य उपसमूह} पढ़िए (अष्टियाँ $\{g : \chi(g) = \chi(e)\}$
और उनके प्रतिच्छेदन)।
\end{exercise}

\begin{exercise}[$\star\star$]\label{exo:b3:representations:6}
(क) दर्शाइए कि $\ker\chi = \{g : \chi(g) = \chi(e)\}$ अंतर्निहित \omterm{def:b3:representations:rep}{निरूपण} की अष्टि है (\cref{prop:b3:representations:charbasics}, समता
स्थिति)।
(ख) दर्शाइए कि $G$ का प्रत्येक \omterm{def:b3:groups:normal}{प्रसामान्य उपसमूह} \omterm{def:b3:representations:rep}{अखंडनीय} वर्णों
की अष्टियों का प्रतिच्छेदन है। \emph{($G/N$ को निष्ठ रूप से
निरूपित कीजिए: उसका \omterm{ex:b3:representations:examples}{नियमित निरूपण}।)}
(ग) निष्कर्ष निकालिए: $G$ सरल है तभी और केवल तभी जब प्रत्येक
अतुच्छ \omterm{def:b3:representations:rep}{अखंडनीय} $\chi_i$ के लिए $\ker\chi_i = \{e\}$ --- अर्थात् सरलता \omterm{def:b3:representations:table}{वर्ण सारणी}
में पढ़ी जा सकती है।
\end{exercise}

\begin{exercise}[$\star\star$]\label{exo:b3:representations:7}
$\abs G = 8$ की घातें: दर्शाइए कि क्रम $8$ के किसी अनाबेली समूह का घात
प्रारूप $(1,1,1,1,2)$ होता है, और उसका घात-$2$ \omterm{def:b3:representations:rep}{निरूपण} निष्ठ होता है।
अधिक व्यापक रूप से दर्शाइए कि क्रम $p^3$ के किसी अनाबेली समूह का
प्रारूप $(1^{\,p^2},
p, \dots, p)$ होता है, जिसमें $p^2$ इकाइयाँ और घात $p$ के
$p - 1$ \omterm{def:b3:representations:character}{वर्ण} होते हैं। \emph{($[G : D(G)]$ और $\sum n_i^2 = \abs G$ का उपयोग कीजिए; यहाँ
$D(G) = Z(G)$ का क्रम $p$ है।)}
\end{exercise}

\begin{exercise}[$\star\star\star$]\label{exo:b3:representations:8}
परिमित समूहों $G, H$ के लिए: दर्शाइए कि $G, H$ के \omterm{def:b3:representations:rep}{अखंडनीय} वर्णों
$\chi, \psi$ के लिए $G \times H$ पर \omterm{def:b3:representations:character}{वर्ग फलन} $\chi(g)\psi(h)$ ठीक $G \times H$ के \omterm{def:b3:representations:rep}{अखंडनीय} \omterm{def:b3:representations:character}{वर्ण}
हैं। \emph{(प्रसामान्य लांबिकता सीधा परिकलन है; और पूर्णता वर्ग
गिनकर।)} इससे $\Z/2\Z \times \Z/2\Z$ की \omterm{def:b3:representations:table}{वर्ण सारणी} निकालिए और संरचना प्रमेय के
माध्यम से परिमित आबेली समूहों के लिए \cref{prop:b3:representations:onedim}(क) पुनः व्युत्पन्न
कीजिए।
\end{exercise}

\begin{exercise}[$\star\star\star$]\label{exo:b3:representations:9}
$A_5$ की \omterm{def:b3:representations:table}{वर्ण सारणी} (\cref{exo:b3:groups:11} से वर्गों के आकार $1, 15, 20, 12,
12$):
(क) दर्शाइए कि घातें $1, 3, 3, 4, 5$ हैं \emph{($n_1 = 1$ वाला $\sum n_i^2 = 60$ का
एकमात्र हल, और सरलता के साथ \cref{exo:b3:representations:6}(ग) का उपयोग करने पर कोई अन्य
$n_i = 1$ नहीं)}।
(ख) घात-$4$ \omterm{def:b3:representations:character}{वर्ण} ($5$ बिंदुओं पर क्रमचय क्रिया) और घात-$5$
\omterm{def:b3:representations:character}{वर्ण} की रचना कीजिए (छह $5$-सिलो उपसमूहों पर क्रिया घात $6 = 1 + 5$
देती है; अखंडनीयता जाँचिए), और दोनों घात-$3$ पंक्तियाँ स्तंभ
लांबिकता से पूरी कीजिए: $5$-चक्र वर्गों पर स्वर्णिम-अनुपात
प्रविष्टियाँ $\frac{1\pm\sqrt5}2$ प्रकट होती हैं।
(ग) पूर्ण सारणी से सत्यापित कीजिए कि $A_5$ सरल है (\cref{exo:b3:representations:6}(ग))।
\end{exercise}

\begin{exercise}[$\star\star$]\label{exo:b3:representations:10}
मान लीजिए $\rho$ घात $n$ का \omterm{def:b3:representations:rep}{अखंडनीय निरूपण} है और $z
\in Z(G)$।
दर्शाइए कि $\rho(z) = \lambda_z\,\mathrm{id}$, जहाँ $\lambda \colon Z(G) \to \C^\times$ एक समाकारिता है (\emph{केंद्रीय
\omterm{def:b3:representations:character}{वर्ण}}), और केंद्रीय $z$ के लिए $\abs{\chi(z)} = n$ निष्कर्ष निकालिए।
अनुप्रयोग: यदि $G$ का कोई निष्ठ \omterm{def:b3:representations:rep}{अखंडनीय निरूपण} हो, तो $Z(G)$
चक्रीय है।
\end{exercise}

\begin{exercise}[$\star\star$]\label{exo:b3:representations:11}
(समजातीय प्रक्षेप) मान लीजिए $(V, \rho)$ $G$ का एक \omterm{def:b3:representations:rep}{निरूपण} है और
$\chi_i$ घात $n_i$ का \omterm{def:b3:representations:rep}{अखंडनीय} \omterm{def:b3:representations:character}{वर्ण}। परिभाषित कीजिए
\[
p_i = \frac{n_i}{\abs G}\sum_{g\in G}
\overline{\chi_i(g)}\,\rho(g) \;\in\; \mathcal L(V).
\]
(क) दर्शाइए कि $p_i$ $G$-समतुल्यवर्ती है, और \omterm{def:b3:representations:character}{वर्ण} $\chi_j$ वाले
किसी \omterm{def:b3:representations:rep}{अखंडनीय} उपनिरूपण $W \subseteq
V$ पर उसका प्रतिबंध परिकलित कीजिए: वह
$\delta_{ij}\,
\mathrm{id}_W$ है \emph{(शूर; अदिश पहचानने के लिए अनुरेख लीजिए)}।
(ख) निष्कर्ष निकालिए कि $p_i$ \omterm{def:b3:representations:character}{वर्ण} $\chi_i$ वाले समस्त \omterm{def:b3:representations:rep}{अखंडनीय}
उपनिरूपणों के योग $V_i$ पर प्रक्षेप है (\emph{समजातीय घटक}), कि
$\sum_ip_i =
\mathrm{id}_V$, और यह कि विघटन $V =
\bigoplus_iV_i$ विहित है --- प्रत्येक $V_i$ के
\omterm{def:b3:representations:rep}{अखंडनीय} निरूपणों में सूक्ष्मतर विभाजन के विपरीत।
(ग) $S_3$ के \omterm{ex:b3:representations:examples}{नियमित निरूपण} और चिह्न \omterm{def:b3:representations:character}{वर्ण} $\varepsilon$ के लिए $p_\varepsilon$ को
समूह बीजगणित के अवयव के रूप में स्पष्ट लिखिए और $p_\varepsilon^2
= p_\varepsilon$ हाथ से
जाँचिए।
\end{exercise}

\begin{exercise}[$\star\star$]\label{exo:b3:representations:12}
(सारणी पढ़ना) क्रम $24$ के किसी समूह $G$ की \omterm{def:b3:representations:table}{वर्ण सारणी} अंशतः
ज्ञात है: उसमें $5$ वर्ग हैं, जिनके आकार $1, 6, 8, 6, 3$ हैं, और घातें
$1, 1, 2, 3, 3$ हैं।
(क) पूरी सारणी पुनः प्राप्त कीजिए: दो रैखिक \omterm{def:b3:representations:character}{वर्ण} (एक तुच्छ;
दूसरा ठीक आकार $6$ और $6$ वाले वर्गों पर मान $-1$ लेता
है), फिर तत्समक स्तंभ के साथ स्तंभ लांबिकता से घात-$2$ \omterm{def:b3:representations:character}{वर्ण},
और उसी प्रकार दोनों घात-$3$ \omterm{def:b3:representations:character}{वर्ण} (उनमें से एक $\chi_2\chi_4$ है)।
(ख) $G$ की पहचान कीजिए ($\cong S_4$: वर्गों की तुलना चक्र प्रकारों
से कीजिए), और \cref{exo:b3:representations:6} के माध्यम से सारणी से \omterm{def:b3:groups:normal}{प्रसामान्य उपसमूह}
निकालिए: $\chi_2$ की अष्टियाँ (सूचकांक $2$: $A_4$) और $\chi_3$ की
(क्लाइन $V_4$), तथा $\{e\}, G$ के अतिरिक्त और कुछ नहीं।
(ग) समझाइए कि सारणी $G/V_4 \cong S_3$ कैसे दिखाती है \emph{(कौन-से \omterm{def:b3:representations:character}{वर्ण}
भागफल से होकर गुजरते हैं?)}।
\end{exercise}

\section{समस्या: बर्नसाइड की \texorpdfstring{$p^aq^b$}{p^a q^b} प्रमेय}

\begin{problem}[{सप्ताहांत समस्या --- क्रम $p^aq^b$ के समूहों की
विलेयता}]\label{pb:b3:representations:1}
बर्नसाइड ने 1904 में सिद्ध किया कि जिस समूह के क्रम में अधिकतम
दो अभाज्य गुणनखंड हों वह \omterm{def:b3:groups:derived}{विलेय} होता है --- अमूर्त समूहों के बारे
में एक ऐसा कथन जिसकी आधी सदी तक ज्ञात सभी उपपत्तियाँ \omterm{def:b3:representations:character}{वर्ण}
सिद्धांत से होकर जाती थीं। यह समस्या \cref{ch:b3:groups} (\omterm{def:b3:groups:derived}{विलेयता}), \cref{ch:b3:modules}
(\omterm{def:b3:modules:free}{परिमिततः जनित} $\Z$-मॉड्यूल) और इस अध्याय को जोड़कर पूरी उपपत्ति
खड़ी करती है। सर्वत्र $\chi_1, \dots, \chi_r$ $G$ के \omterm{def:b3:representations:rep}{अखंडनीय} \omterm{def:b3:representations:character}{वर्ण} हैं और $n_i = \chi_i(e)$।

\medskip
\noindent\textbf{भाग I --- \omterm{def:b3:galois:algebraic}{बीजीय} पूर्णांक।} \emph{\omterm{def:b3:galois:algebraic}{बीजीय}
पूर्णांक} $\Z[X]$ के किसी \emph{मोनिक} बहुपद का मूल होता है।
\begin{enumerate}
  \item दर्शाइए कि $\alpha$ \omterm{def:b3:galois:algebraic}{बीजीय} पूर्णांक है तभी और केवल तभी जब
        वलय $\Z[\alpha]$ एक \omterm{def:b3:modules:free}{परिमिततः जनित} $\Z$-मॉड्यूल हो।
  \item निष्कर्ष निकालिए कि \omterm{def:b3:galois:algebraic}{बीजीय} पूर्णांक $\C$ का एक उपवलय
        बनाते हैं। \emph{(यदि $\Z[\alpha], \Z[\beta]$ \omterm{def:b3:modules:free}{परिमिततः जनित} हों, तो $\Z[\alpha, \beta]$
        भी है, और \cref{thm:b3:modules:submodule} तथा एक प्रस्तुति तर्क से \omterm{def:b3:modules:free}{परिमिततः जनित}
        $\Z$-मॉड्यूलों के उपमॉड्यूल \omterm{def:b3:modules:free}{परिमिततः जनित} होते हैं ---
        अथवा सीधे: $\Z^n$ का उपमॉड्यूल कोटि $\leq n$ का \omterm{def:b3:modules:free}{मुक्त
        मॉड्यूल} है।)}
  \item दर्शाइए कि कोई \emph{परिमेय} \omterm{def:b3:galois:algebraic}{बीजीय} पूर्णांक पूर्णांक ही
        होता है। \emph{(परिमेय मूल प्रमेय।)}
  \item दर्शाइए कि प्रत्येक \omterm{def:b3:representations:character}{वर्ण} मान $\chi(g)$ एक \omterm{def:b3:galois:algebraic}{बीजीय} पूर्णांक है।
\end{enumerate}

\medskip
\noindent\textbf{भाग II --- वर्ग-योग संबंध।} घात $n$ और \omterm{def:b3:representations:character}{वर्ण}
$\chi$ वाला कोई \omterm{def:b3:representations:rep}{अखंडनीय} $(V, \rho)$ स्थिर कीजिए। किसी \omterm{ex:b3:groups:actions}{संयुग्मन वर्ग}
$C$ के लिए $S_C = \sum_{g \in C}\rho(g) \in
\mathcal L(V)$ रखिए।
\begin{enumerate}[resume]
  \item दर्शाइए कि $S_C$ समतुल्यवर्ती है, अतः $S_C =
        \omega_C\,\mathrm{id}$, जहाँ
        \[
        \omega_C = \frac{\abs C\,\chi(g_C)}{n}
        \qquad (g_C \in C \text{ कोई भी प्रतिनिधि}).
        \]
  \item दर्शाइए कि $S_CS_{C'} = \sum_{C''}
        a_{CC'C''}\,S_{C''}$, जहाँ $a_{CC'C''} \in \N$ किसी स्थिर $z \in C''$ के लिए
        $xy = z$ वाले युग्मों $(x, y) \in C \times
        C'$ की गणना करता है। इससे निष्कर्ष
        निकालिए कि $\omega_C$ $\omega_C\,\omega_{C'} = \sum_{C''}
        a_{CC'C''}\,\omega_{C''}$ को संतुष्ट करते हैं।
  \item निष्कर्ष निकालिए कि प्रत्येक $\omega_C$ एक \omterm{def:b3:galois:algebraic}{बीजीय} पूर्णांक है।
        \emph{($1$ और $\omega_C$ से जनित $\Z$-मॉड्यूल एक \omterm{def:b3:modules:free}{परिमिततः
        जनित} वलय है; प्रश्न 1 की कसौटी लगाइए --- अधिक ठीक-ठीक,
        दर्शाइए कि $M = \Z\text{-विस्तार}
        (1, (\omega_C)_C)$ $\omega_{C_0} M \subseteq
        M$ को संतुष्ट करता है और
        सारणिक/केली--हैमिल्टन युक्ति का उपयोग कीजिए, अथवा प्रश्न
        2 का उपवलय तर्क।)}
  \item \emph{\omterm{thm:b3:galois:finitefields}{फ्रोबेनियस} विभाज्यता} निकालिए: प्रत्येक \omterm{def:b3:representations:rep}{अखंडनीय} घात
        के लिए $n_i$ $\abs G$ को विभाजित करता है। \emph{($\frac{\abs G}{n} = \frac{\abs G}{n}
        \langle\chi,\chi\rangle = \sum_C \omega_C\,
        \overline{\chi(g_C)}$
        परिकलित कीजिए: एक \omterm{def:b3:galois:algebraic}{बीजीय} पूर्णांक जो परिमेय है।)}
\end{enumerate}

\medskip
\noindent\textbf{भाग III --- बर्नसाइड की सरलता कसौटी।}
\begin{enumerate}[resume]
  \item मान लीजिए $\chi$ घात $n$ का \omterm{def:b3:representations:rep}{अखंडनीय} है और $C$ ऐसा वर्ग
        जिसके लिए $\gcd(\abs C, n) = 1$। बेज़ू और प्रश्न 4--7 का उपयोग करके
        दर्शाइए कि $\frac{\chi(g_C)}{n}$ एक \omterm{def:b3:galois:algebraic}{बीजीय} पूर्णांक है।
  \item अब मान लीजिए इसके अतिरिक्त $0 < \abs{\chi(g_C)} < n$। दर्शाइए कि यह असंभव है:
        \omterm{def:b3:galois:algebraic}{बीजीय} पूर्णांक $\alpha =
        \chi(g_C)/n$ के समस्त \emph{संयुग्म} --- अर्थात्
        $\sigma \in \operatorname{Gal}(\Q(\zeta_{\abs G})/\Q)$ के लिए संख्याएँ $\alpha_\sigma = \frac1n\sigma(\chi(g_C))$, जिनमें से प्रत्येक $n$
        इकाई-मूलों का औसत है --- मापांक $\leq 1$ के हैं, अतः गुणनफल
        $N = \prod_\sigma\alpha_\sigma$ ऐसा परिमेय \omterm{def:b3:galois:algebraic}{बीजीय} पूर्णांक है जिसके लिए $0 < \abs N < 1$ ---
        \omterm{def:b3:galois:galois}{गाल्वा समूह} के लिए \cref{thm:b3:galois:cyclotomicirred} और इस तथ्य को उद्धृत करते हुए कि
        $\sigma$ इकाई-मूलों का क्रमचय करता है, प्रत्येक कथन को उचित
        ठहराइए। निष्कर्ष निकालिए: या तो $\chi(g_C) = 0$ या $\rho(g_C)$ अदिश है
        (\cref{prop:b3:representations:charbasics})।
  \item (बर्नसाइड की कसौटी) मान लीजिए $C \neq \{e\}$ \emph{अभाज्य घात}
        आकार $p^k > 1$ का \omterm{ex:b3:groups:actions}{संयुग्मन वर्ग} है, और मान लीजिए $G$ सरल
        अनाबेली है। $C$ के स्तंभ की $e$ के स्तंभ के साथ स्तंभ
        लांबिकता से $1 + \sum_{i \geq 2} n_i\chi_i(g_C) = 0$ मिलता है। दर्शाइए कि $p \nmid n_i$ वाला कोई
        अतुच्छ $\chi_i$ $\chi_i(g_C) \neq 0$ को संतुष्ट करता है \emph{(अन्यथा
        $\frac1p$ एक \omterm{def:b3:galois:algebraic}{बीजीय} पूर्णांक होता)}; प्रश्न 10 से $\rho_i(g_C)$ अदिश
        है; अब सरलता के साथ विरोधाभास निकालिए \emph{($\rho_i(g)$ के
        अदिश होने वाले $g$ का समुच्चय एक \omterm{def:b3:groups:normal}{प्रसामान्य उपसमूह} है;
        \cref{exo:b3:representations:6} से निष्ठता का उपयोग कीजिए)}। निष्कर्ष निकालिए:
        \emph{किसी भी सरल अनाबेली समूह में $> 1$ अभाज्य घात आकार
        का कोई \omterm{ex:b3:groups:actions}{संयुग्मन वर्ग} नहीं होता।}
\end{enumerate}

\medskip
\noindent\textbf{भाग IV --- प्रमेय।}
\begin{enumerate}[resume]
  \item मान लीजिए $a + b \geq 1$ के साथ $\abs G = p^aq^b$। यदि $G$ सरल हो, तो
        दर्शाइए कि वह आबेली है: किसी सिलो $q$-उपसमूह के \omterm{ex:b3:groups:actions}{केंद्र}
        में $z \neq e$ चुनिए (\cref{thm:b3:groups:pfixed}) और उसके \omterm{ex:b3:groups:actions}{संयुग्मन वर्ग} $[G : Z_G(z)]$ के
        आकार पर विचार कीजिए, जो $p$ की एक घात है (क्यों?);
        अब प्रश्न 11 लगाइए।
  \item $\abs G$ पर आगमन से निष्कर्ष निकालिए: \emph{क्रम $p^aq^b$ का
        प्रत्येक समूह \omterm{def:b3:groups:derived}{विलेय} है} (बर्नसाइड)। तीन अभाज्यों के लिए
        यह तर्क क्यों टूट जाता है --- और $\abs{A_5} = 2^2\cdot3\cdot5$ को देखते हुए उसे
        टूटना ही क्यों चाहिए?
\end{enumerate}

\medskip
\noindent\textbf{भाग V --- $A_5$ की \omterm{def:b3:representations:table}{वर्ण सारणी}।} जिस सबसे छोटे
समूह तक बर्नसाइड की प्रमेय नहीं पहुँच सकती, वह अपने पूरे चित्र का
अधिकारी है; नीचे सब कुछ केवल इस अध्याय और \cref{exo:b3:groups:11} का उपयोग करता
है।
\begin{enumerate}[resume]
  \item \cref{exo:b3:groups:11} से $A_5$ के पाँच \omterm{ex:b3:groups:actions}{संयुग्मन वर्ग} स्मरण कीजिए: $\{e\}$,
        $15$ द्विक स्थानांतरण, $20$ त्रि-चक्र, और प्रत्येक में $12$
        पाँच-चक्र वाले दो वर्ग, जिनके प्रतिनिधि $c =
        (1\,2\,3\,4\,5)$ और $c^2$
        हैं। समझाइए कि पाँच-चक्र $S_5$-वर्ग में एक ही वर्ग बनाते
        हुए भी दो $A_5$-वर्गों में क्यों बँट जाते हैं।
  \item दर्शाइए कि $A_5$ की \omterm{def:b3:representations:rep}{अखंडनीय} घातें ठीक $1, 3, 3, 4, 5$ हैं: $r = 5$
        वर्गों के साथ $\sum_in_i^2 = 60$ का उपयोग कीजिए, और इस तथ्य का कि
        $A_5$ पूर्ण है ($D(A_5) =
        A_5$), अतः तुच्छ \omterm{def:b3:representations:character}{वर्ण} ही उसका एकमात्र
        रैखिक \omterm{def:b3:representations:character}{वर्ण} है; फिर हर दूसरे बहुसमुच्चय को अपवर्जित कीजिए
        \emph{($59$ को $\geq 2$ पूर्णांकों के चार वर्गों के योग के
        रूप में लिखिए: जाँचिए कि यह केवल एक ही प्रकार से होता है
        जिसमें सभी योज्य विश्वसनीय घातें हों)}।
  \item मान लीजिए $\pi$ $\{1, \dots, 5\}$ पर $A_5$ का क्रमचय \omterm{def:b3:representations:character}{वर्ण} है:
        $\pi(g) = \#\operatorname{Fix}(g)$, जिसके पाँचों वर्गों पर मान $5, 1, 2, 0, 0$ हैं। $\langle\pi, \mathbf 1\rangle$ और
        $\langle\pi,
        \pi\rangle$ परिकलित कीजिए, और निष्कर्ष निकालिए कि $\chi_4 = \pi - \mathbf 1$ घात
        $4$ का \omterm{def:b3:representations:rep}{अखंडनीय} है, जिसके मान $4, 0, 1, -1,
        -1$ हैं।
  \item $10$ अक्रमित युग्मों $\{i, j\}$ पर वही खेल: स्थिर-बिंदु
        गणनाएँ $10, 2, 1, 0, 0$ हैं। $\langle\pi_{10}, \pi_{10}\rangle$, $\langle\pi_{10},
        \mathbf 1\rangle$ और $\langle\pi_{10}, \chi_4\rangle$ परिकलित कीजिए,
        विघटन $\pi_{10} = \mathbf 1 + \chi_4
        + \chi_5$ निकालिए, और घात $5$ का \omterm{def:b3:representations:rep}{अखंडनीय} $\chi_5$
        प्राप्त कीजिए, जिसके मान $5, 1, -1, 0, 0$ हैं।
  \item शेष दो \omterm{def:b3:representations:rep}{अखंडनीय} $\chi_2, \chi_3$ की घात $3$ है। स्तंभ लांबिकता
        (प्रत्येक अतत्समक स्तंभ की तत्समक स्तंभ के साथ, और
        प्रत्येक स्तंभ की अपने साथ) पाँच-चक्रों के बाहर उनके मान
        निर्धारित कर देती है: दर्शाइए कि द्विक स्थानांतरणों पर
        $\chi_2(g) = \chi_3(g) = -1$ और त्रि-चक्रों पर $0$।
  \item पाँच-चक्र वर्गों पर $x = \chi_2(c)$ और $y =
        \chi_2(c^2)$ रखिए; सममिति से
        $\chi_3(c) = y$, $\chi_3(c^2) = x$ लिया जा सकता है। $c$ के स्तंभ को तत्समक
        स्तंभ के साथ और $c^2$ के स्तंभ के साथ युग्मित करके $x + y = 1$
        और $xy = -1$ निकालिए (और $c$ के स्तंभ की अपने साथ युग्मन से
        मिलने वाले मान $x^2 +
        y^2 = 3$ की जाँच कीजिए), जिससे
        \[
        \{x, y\} = \Bigl\{\frac{1 + \sqrt5}2,\ \frac{1 -
        \sqrt5}2\Bigr\} :
        \]
        अर्थात् स्वर्णिम अनुपात और उसका संयुग्म। अब $A_5$ की पूरी
        \omterm{def:b3:representations:table}{वर्ण सारणी} संकलित कीजिए।
  \item जाँचें चलाइए: $\chi_2$ की पंक्ति मानक $1$ है ($\varphi^2 + \bar\varphi^2 = 3$ का
        उपयोग कीजिए), $\langle\chi_2,
        \chi_3\rangle = 0$, और पाँचों घातों के लिए \omterm{thm:b3:galois:finitefields}{फ्रोबेनियस}
        विभाज्यता (प्रश्न 8)। सारणी में कहाँ आप $S_5$ से अंतर
        \emph{देखते} हैं, जिसके सभी \omterm{def:b3:representations:character}{वर्ण} मान परिमेय पूर्णांक हैं?
  \item केवल सारणी से निष्कर्ष निकालिए कि $A_5$ सरल है: कोई
        \omterm{def:b3:groups:normal}{प्रसामान्य उपसमूह} $e$ को अंतर्विष्ट करने वाले संयुग्मन
        वर्गों का सम्मिलन होता है जिसकी गणनीयता $60$ को विभाजित
        करती है --- जाँचिए कि $1 + 15 + 20 + 12 + 12$ का कोई भी ऐसा उचित उपयोग जिसमें
        पद $1$ हो, $60$ को विभाजित नहीं करता। प्रश्न 11 की
        कसौटी से मिलान कीजिए: सत्यापित कीजिए कि $A_5$ के किसी वर्ग
        का आकार अभाज्य घात $> 1$ नहीं है।
  \item (बीसफलकी समापन) $A_5$ बीसफलक का घूर्णन समूह है, और घात-$3$
        \omterm{def:b3:representations:rep}{निरूपण} $\R^3$ पर दोनों ज्यामितीय क्रियाएँ हैं। अनुरेख
        सर्वसमिका सत्यापित कीजिए: कोण $\theta$ के घूर्णन का अनुरेख
        $1 +
        2\cos\theta$ होता है, और $1 + 2\cos\frac{2\pi}5 =
        \frac{1+\sqrt5}2 = \varphi$। बिना किसी परिकलन के समझाइए कि
        दूसरे घात-$3$ \omterm{def:b3:representations:character}{वर्ण} को संयुग्म मान क्यों धारण करना चाहिए:
        $\Q(\sqrt5)/\Q$ का \omterm{def:b3:galois:galois}{गाल्वा समूह} पूरी \omterm{def:b3:representations:table}{वर्ण सारणी} पर (प्रविष्टिशः)
        क्रिया करता है और \omterm{def:b3:representations:rep}{अखंडनीय} वर्णों का क्रमचय कर देता है।
\end{enumerate}

\medskip
\noindent\textbf{भाग VI --- पूरक: एक केंद्रीय परिबंध और टेंसर
वर्ग।}
\begin{enumerate}[resume]
  \item (विभाज्यता से तीक्ष्णतर) मान लीजिए $\chi$ घात $n$ का
        \omterm{def:b3:representations:rep}{अखंडनीय} है। दर्शाइए कि प्रत्येक $z \in Z(G)$ के लिए $\abs{\chi(z)} = n$
        \emph{(शूर की प्रमेयिका: $\rho(z)$ परिमित क्रम का अदिश है)},
        और $\langle\chi, \chi\rangle = 1$ से परिबंध
        \[
        n^2 \leq [G : Z(G)] .
        \]
        निकालिए। दर्शाइए कि क्रम $8$ के अनाबेली समूहों की
        \omterm{def:b3:representations:rep}{अखंडनीय} घातें $1, 1, 1, 1, 2$ हैं \emph{(पाँच \omterm{ex:b3:groups:actions}{संयुग्मन वर्ग};
        $8$ को पाँच वर्गों के योग के रूप में लिखिए)} और वे समता
        $4 = [G : Z(G)]$ प्राप्त कर लेते हैं; $A_5$ पर परिबंध जाँचिए, जिसका
        \omterm{ex:b3:groups:actions}{केंद्र} तुच्छ है।
  \item ($\chi_2$ का टेंसर वर्ग) परिमित क्रम के $g$ के लिए $\rho(g)$
        इकाई-मूल अभिलक्षणिक मानों के साथ विकर्णनीय है; इससे \omterm{def:b3:representations:character}{वर्ण}
        सूत्र
        \[
        \chi_{\operatorname{Sym}^2 V}(g)
        = \frac{\chi(g)^2 + \chi(g^2)}2,
        \qquad
        \chi_{\Lambda^2 V}(g)
        = \frac{\chi(g)^2 - \chi(g^2)}2 .
        \]
        निकालिए। इन्हें $A_5$ के $\chi_2$ पर लगाइए \emph{(ध्यान
        दीजिए कि जब $g$ $c$ के वर्ग में विचरता है तब $g^2$
        $c^2$ के वर्ग में विचरता है, और विलोमतः भी)}: दर्शाइए कि
        $\Lambda^2\chi_2 =
        \chi_2$ और $\operatorname{Sym}^2\chi_2 = \mathbf 1 +
        \chi_5$, अतः
        \[
        \chi_2\otimes\chi_2 = \mathbf 1 + \chi_2 + \chi_5 .
        \]
        $\R^3$ पर सदिश गुणन के माध्यम से $\Lambda^2\chi_2 = \chi_2$ की ज्यामितीय
        व्याख्या कीजिए।
  \item (सारणी का अंतिम अंकेक्षण) संख्यात्मक रूप से सत्यापित
        कीजिए: दोनों पाँच-चक्र स्तंभों के बीच स्तंभ लांबिकता
        ($\varphi\bar\varphi = -1$), $c$ के स्तंभ की अपने साथ युग्मन से मिलने वाला
        मान $5 =
        \abs{Z_{A_5}(c)}$, और सारणी के चारों अतत्समक स्तंभों पर नियमित
        \omterm{def:b3:representations:character}{वर्ण} $\sum_i
        n_i\chi_i(g) = 0$ का लुप्त होना।
\end{enumerate}
\end{problem}
